% GENERATED FILE. Edit canonical lessons, then run npm run generate:books.
% canonical-source-hash: fnv1a64:bd4681669c28e617
% canonical-lessons: UR-C63-javab, UR-C63-lafz, UR-C63-khel, UR-C63-darakht, UR-C63-phul

\chapter{Answer, Word, Game, Tree, Flower}
\label{ch:ur-answer-word-game-tree-flower}
\hlchaptermodality{63}

\begin{chapteropening}
\textbf{By the end of this chapter:} \emph{I can say an answer, a word, a game, a tree and a flower.}

Complete the last lesson of chapter 63: I can say an answer, a word, a game, a tree and a flower.
\end{chapteropening}

\section[\texorpdfstring{\ur{جواب} (javāb)}{جواب (javāb)}]{\ur{جواب} (javāb) — an answer}

\label{lesson:UR-C63-javab}



\begin{quote}
Before the new one: say the Urdu for a lesson, then the Urdu for a question.
\end{quote}

\subsection*{You'll want to know: \ur{جواب}}
\textbf{\ur{جواب}} — \emph{javāb} — \textquotedblleft{}an answer\textquotedblright{}.

\begin{cousinweb}
\textbf{\ur{جواب}} (\emph{javāb}) is a loan from Arabic \emph{jawāb}, \textquotedblleft{}an answer\textquotedblright{}, which reached Urdu through Persian. Its Arabic root is \emph{j-w-b} (Lane).
\end{cousinweb}

\begin{grammarlens}[title={What you've built}]
One more word for this chapter.
\end{grammarlens}

\subsection*{Your turn}
\begin{itemize}
\raggedright
  \item \emph{Say it:} \emph{javāb}
  \item \emph{Say it:} \emph{javāb}, once more
  \item \emph{Say it:} say \emph{javāb}
  \item \emph{Recall:} say the Urdu for a lesson, then the Urdu for a question, then say \emph{javāb} again
\end{itemize}

\begin{tcolorbox}[breakable,colback=teal!4,colframe=teal!35!black,title={Before you move on}]
What does \ur{جواب} mean? (\textquotedblleft{}An answer\textquotedblright{}.) Say it once more.
\end{tcolorbox}
\section[\texorpdfstring{\ur{لفظ} (lafz)}{لفظ (lafz)}]{\ur{لفظ} (lafz) — a word}

\label{lesson:UR-C63-lafz}



\begin{quote}
Before the new one: say the Urdu for a question, then the Urdu for an answer.
\end{quote}

\subsection*{You'll want to know: \ur{لفظ}}
\textbf{\ur{لفظ}} — \emph{lafz} — \textquotedblleft{}a word\textquotedblright{}.

\begin{cousinweb}
\textbf{\ur{لفظ}} (\emph{lafz}) is a loan from Arabic \emph{lafẓ}, \textquotedblleft{}utterance, word\textquotedblright{}, which reached Urdu through Persian. Its Arabic root is \emph{l-f-ẓ} (Lane).
\end{cousinweb}

\begin{grammarlens}[title={What you've built}]
One more word for this chapter.
\end{grammarlens}

\subsection*{Your turn}
\begin{itemize}
\raggedright
  \item \emph{Say it:} \emph{lafz}
  \item \emph{Say it:} \emph{lafz}, once more
  \item \emph{Say it:} say \emph{lafz}
  \item \emph{Recall:} say the Urdu for a question, then the Urdu for an answer, then say \emph{lafz} again
\end{itemize}

\begin{tcolorbox}[breakable,colback=teal!4,colframe=teal!35!black,title={Before you move on}]
What does \ur{لفظ} mean? (\textquotedblleft{}A word\textquotedblright{}.) Say it once more.
\end{tcolorbox}
\section[\texorpdfstring{\ur{کھیل} (khel)}{کھیل (khel)}]{\ur{کھیل} (khel) — a game}

\label{lesson:UR-C63-khel}



\begin{quote}
Before the new one: say the Urdu for an answer, then the Urdu for a word.
\end{quote}

\subsection*{You'll want to know: \ur{کھیل}}
\textbf{\ur{کھیل}} — \emph{khel} — \textquotedblleft{}a game\textquotedblright{}.

\begin{cousinweb}
\textbf{\ur{کھیل}} (\emph{khel}) is inherited from a reconstructed \emph{*khel-}, \textquotedblleft{}play\textquotedblright{} (CDIAL 3918). CDIAL links it with Sanskrit \emph{krīḍati}, \textquotedblleft{}plays\textquotedblright{}: if that verb has a non-Aryan source, these forms may share it.
\end{cousinweb}

\begin{grammarlens}[title={What you've built}]
One more word for this chapter.
\end{grammarlens}

\subsection*{Your turn}
\begin{itemize}
\raggedright
  \item \emph{Say it:} \emph{khel}
  \item \emph{Say it:} \emph{khel}, once more
  \item \emph{Say it:} say \emph{khel}
  \item \emph{Recall:} say the Urdu for an answer, then the Urdu for a word, then say \emph{khel} again
\end{itemize}

\begin{tcolorbox}[breakable,colback=teal!4,colframe=teal!35!black,title={Before you move on}]
What does \ur{کھیل} mean? (\textquotedblleft{}A game\textquotedblright{}.) Say it once more.
\end{tcolorbox}
\section[\texorpdfstring{\ur{درخت} (darakht)}{درخت (darakht)}]{\ur{درخت} (darakht) — a tree}

\label{lesson:UR-C63-darakht}



\begin{quote}
Before the new one: say the Urdu for a word, then the Urdu for a game.
\end{quote}

\subsection*{You'll want to know: \ur{درخت}}
\textbf{\ur{درخت}} — \emph{darakht} — \textquotedblleft{}a tree\textquotedblright{}.

\begin{cousinweb}
\textbf{\ur{درخت}} (\emph{darakht}) is a loan, usually traced to Persian \emph{darakht}, \textquotedblleft{}tree\textquotedblright{}.
\end{cousinweb}

\begin{grammarlens}[title={What you've built}]
One more word for this chapter.
\end{grammarlens}

\subsection*{Your turn}
\begin{itemize}
\raggedright
  \item \emph{Say it:} \emph{darakht}
  \item \emph{Say it:} \emph{darakht}, once more
  \item \emph{Say it:} say \emph{darakht}
  \item \emph{Recall:} say the Urdu for a word, then the Urdu for a game, then say \emph{darakht} again
\end{itemize}

\begin{tcolorbox}[breakable,colback=teal!4,colframe=teal!35!black,title={Before you move on}]
What does \ur{درخت} mean? (\textquotedblleft{}A tree\textquotedblright{}.) Say it once more.
\end{tcolorbox}
\section[\texorpdfstring{\ur{پھول} (phūl)}{پھول (phūl)}]{\ur{پھول} (phūl) — a flower}

\label{lesson:UR-C63-phul}



\begin{quote}
Before the new one: say the Urdu for a game, then the Urdu for a tree.
\end{quote}

\subsection*{You'll want to know: \ur{پھول}}
\textbf{\ur{پھول}} — \emph{phūl} — \textquotedblleft{}a flower\textquotedblright{}.

\begin{cousinweb}
\textbf{\ur{پھول}} (\emph{phūl}) is a tadbhava, a word inherited from Sanskrit and reshaped by sound change. It comes from \emph{phulla}, \textquotedblleft{}blown (of flowers)\textquotedblright{} (CDIAL 9092).
\end{cousinweb}

\begin{grammarlens}[title={What you've built}]
One more word for this chapter.
\end{grammarlens}

\subsection*{Your turn}
\begin{itemize}
\raggedright
  \item \emph{Say it:} \emph{phūl}
  \item \emph{Say it:} \emph{phūl}, once more
  \item \emph{Say it:} say \emph{phūl}
  \item \emph{Recall:} say the Urdu for a game, then the Urdu for a tree, then say \emph{phūl} again
\end{itemize}

\begin{tcolorbox}[breakable,colback=teal!4,colframe=teal!35!black,title={Before you move on}]
What does \ur{پھول} mean? (\textquotedblleft{}A flower\textquotedblright{}.) Say it once more.
\end{tcolorbox}
